% autenticacion_seguridad.tex — Capítulo 6: Autenticación y Seguridad
\chapter{Autenticación y Seguridad}

\section{Estrategia de Autenticación}

FINARA implementa autenticación \textbf{stateless} mediante \textbf{JSON Web Tokens (JWT)}, estándar definido en el RFC 7519 \cite{jwt_rfc}. Esta estrategia elimina la necesidad de almacenar sesiones en el servidor, simplificando el escalado horizontal y la comunicación con la app móvil.

\subsection{JSON Web Token (JWT)}

El token generado al hacer login contiene tres secciones codificadas en Base64URL separadas por puntos:

\begin{lstlisting}[style=json, caption={Estructura del JWT generado por FINARA}]
// Header
{ "alg": "HS256", "typ": "JWT" }

// Payload (claims)
{
  "userId": 1,
  "email":  "juan@email.com",
  "iat":    1741000000,   // issued at (Unix timestamp)
  "exp":    1741604800    // expires at (+7 dias)
}

// Firma: HMACSHA256(base64(header) + "." + base64(payload), JWT_SECRET)
\end{lstlisting}

\begin{lstlisting}[style=javascript, caption={Generación del token JWT en authController.js}]
const jwt = require('jsonwebtoken');

// Al hacer login o register exitoso:
const token = jwt.sign(
  { userId: user.id_usuario, email: user.email },
  process.env.JWT_SECRET,         // Almacenado en .env, nunca en codigo
  { expiresIn: process.env.JWT_EXPIRES_IN || '7d' }
);
\end{lstlisting}

\subsection{Hash de Contraseñas con bcrypt}

Las contraseñas se almacenan exclusivamente como hash bcrypt con \textbf{10 rounds de sal} \cite{bcrypt_npm}. El factor de costo 10 implica $2^{10} = 1024$ iteraciones del algoritmo, haciendo impráctico el ataque por fuerza bruta.

\begin{lstlisting}[style=javascript, caption={Hash y verificación de contraseña con bcrypt}]
const bcrypt = require('bcrypt');
const SALT_ROUNDS = 10;

// Al registrar usuario:
const password_hash = await bcrypt.hash(password, SALT_ROUNDS);

// Al hacer login:
const esValida = await bcrypt.compare(passwordIngresada, user.password_hash);
if (!esValida) return errorResponse(res, 'Credenciales incorrectas', 401);
\end{lstlisting}

\section{Middleware de Autenticación}

Todas las rutas protegidas pasan por \texttt{src/middleware/auth.js} antes de llegar al controlador.

\begin{lstlisting}[style=javascript, caption={Middleware de autenticación JWT completo}]
// src/middleware/auth.js
const jwt = require('jsonwebtoken');

const authMiddleware = (req, res, next) => {
  try {
    const authHeader = req.headers['authorization'];

    if (!authHeader || !authHeader.startsWith('Bearer ')) {
      return res.status(401).json({
        success: false,
        error: 'Token de autenticacion requerido',
      });
    }

    const token   = authHeader.split(' ')[1];
    const decoded = jwt.verify(token, process.env.JWT_SECRET);

    // Adjunta datos del usuario al objeto request
    req.user = { userId: decoded.userId, email: decoded.email };

    next();
  } catch (error) {
    if (error.name === 'TokenExpiredError') {
      return res.status(401).json({
        success: false,
        error: 'Token expirado, inicia sesion nuevamente',
      });
    }
    return res.status(401).json({ success: false, error: 'Token invalido' });
  }
};

module.exports = authMiddleware;
\end{lstlisting}

\section{Validación de Inputs con express-validator}

Cada endpoint con datos de entrada implementa validación mediante \texttt{express-validator} \cite{express_validator}. Las validaciones se declaran en la definición de rutas y son procesadas por el middleware centralizado \texttt{validator.js}.

\begin{lstlisting}[style=javascript, caption={Validaciones del endpoint de registro}]
// src/routes/auth.js
const { body } = require('express-validator');

router.post('/register',
  [
    body('nombre')
      .trim()
      .notEmpty().withMessage('El nombre es requerido')
      .isLength({ min: 2 }).withMessage('Minimo 2 caracteres'),
    body('email')
      .isEmail().withMessage('Email invalido')
      .normalizeEmail(),
    body('password')
      .isLength({ min: 6 }).withMessage('La contrasena debe tener minimo 6 caracteres'),
  ],
  validate,        // middleware centralizado
  register         // controlador
);
\end{lstlisting}

\section{Sanitización de Inputs}

El middleware \texttt{sanitizer.js} escapa caracteres HTML en todos los campos de texto del body antes de que lleguen al controlador, previniendo ataques XSS \cite{owasp_xss}.

\begin{lstlisting}[style=javascript, caption={Middleware sanitizador de inputs}]
// src/middleware/sanitizer.js
const sanitizeString = (str) => {
  if (typeof str !== 'string') return str;
  return str
    .replace(/</g, '&lt;').replace(/>/g, '&gt;')
    .replace(/"/g, '&quot;').replace(/'/g, '&#x27;')
    .trim();
};

const sanitizer = (req, res, next) => {
  if (req.body && typeof req.body === 'object') {
    req.body = sanitizeObject(req.body);
  }
  next();
};
\end{lstlisting}

\section{Rate Limiting}

El middleware \texttt{rateLimiter.js} implementa limitación de tasa de solicitudes para prevenir ataques de fuerza bruta y abuso de la API \cite{owasp_brute_force}. Usa un Map en memoria (adecuado para una sola instancia en desarrollo).

\begin{longtable}{>{\ttfamily}p{3.5cm} >{\centering}p{3cm} >{\centering}p{3cm} p{4.5cm}}
  \caption{Configuración de rate limiting por contexto}
  \label{tab:rate-limiting} \\
  \toprule
  \textbf{Contexto} & \textbf{Máx. Solicitudes} & \textbf{Ventana} & \textbf{Justificación} \\
  \midrule
  \texttt{/api/auth/*}     & 10  & 15 min & Previene fuerza bruta en login \\
  \texttt{/api/tickets/scan} & 20  & 15 min & OCR es costoso en CPU/tiempo \\
  \texttt{/api/*} (general) & 200 & 15 min & Uso normal de la API \\
  \bottomrule
\end{longtable}

\section{Headers de Seguridad}

En lugar de añadir la dependencia \texttt{helmet}, los headers de seguridad se configuran manualmente en \texttt{app.js} \cite{owasp_headers}:

\begin{lstlisting}[style=javascript, caption={Headers de seguridad HTTP configurados manualmente}]
app.use((req, res, next) => {
  res.setHeader('X-Content-Type-Options', 'nosniff');
  // Previene MIME sniffing

  res.setHeader('X-Frame-Options', 'DENY');
  // Previene clickjacking (no iframes)

  res.setHeader('X-XSS-Protection', '1; mode=block');
  // Activa filtro XSS del navegador

  res.setHeader('Referrer-Policy', 'no-referrer');
  // No envia Referer en solicitudes cross-origin

  next();
});
\end{lstlisting}

\section{Prevención de Inyección SQL}

El driver \texttt{mysql2} utiliza \textbf{prepared statements} de forma nativa, enviando los parámetros separados de la consulta SQL. Esto hace imposible la inyección SQL independientemente del valor del input \cite{owasp_sqli}.

\begin{lstlisting}[style=javascript, caption={Prepared statement con mysql2}]
// Correcto: parametro separado de la consulta
const [rows] = await pool.query(
  'SELECT * FROM Usuarios WHERE email = ?',
  [email]   // <- mysql2 envia esto como parametro, nunca interpolado
);

// Incorrecto (no se hace en FINARA):
// pool.query(`SELECT * FROM Usuarios WHERE email = '${email}'`)
\end{lstlisting}

\section{Resumen de Medidas de Seguridad}

\begin{longtable}{p{4.5cm} p{4cm} p{5cm}}
  \caption{Resumen de medidas de seguridad implementadas}
  \label{tab:seguridad-resumen} \\
  \toprule
  \textbf{Amenaza} & \textbf{Mitigación} & \textbf{Implementación} \\
  \midrule
  Acceso no autorizado & JWT + expiración 7d & \texttt{middleware/auth.js} \\
  Contraseñas robadas  & bcrypt 10 rounds    & \texttt{authController.js} \\
  Fuerza bruta login   & Rate limit 10/15min & \texttt{middleware/rateLimiter.js} \\
  XSS en inputs        & Escapado HTML       & \texttt{middleware/sanitizer.js} \\
  Inyección SQL        & Prepared statements & \texttt{mysql2} driver \\
  Clickjacking         & X-Frame-Options     & \texttt{app.js} headers \\
  MIME sniffing        & X-Content-Type      & \texttt{app.js} headers \\
  Archivos maliciosos  & Multer tipo+tamaño  & \texttt{config/upload.js} \\
  \bottomrule
\end{longtable}
